\documentclass[UTF8,11pt]{ctexart}
\usepackage[a4paper,margin=24mm]{geometry}
\usepackage{amsmath,amssymb,amsthm,mathtools,mathrsfs}
\usepackage[all]{xy}
\usepackage{xurl,hyperref,enumitem}
\usepackage{tikz}
\usetikzlibrary{arrows.meta,cd,decorations.markings,calc,patterns}
\hypersetup{hidelinks,breaklinks=true}
\setlength{\emergencystretch}{3em}
\allowdisplaybreaks
\newtheorem*{theorem}{Theorem}
\newtheorem*{thm}{Theorem}
\newtheorem*{lemma}{Lemma}
\newtheorem*{proposition}{Proposition}
\newtheorem*{prop}{Proposition}
\newtheorem*{corollary}{Corollary}
\newtheorem*{cor}{Corollary}
\newtheorem*{claim}{Claim}
\theoremstyle{definition}
\newtheorem*{definition}{Definition}
\newtheorem*{defn}{Definition}
\newtheorem*{remark}{Remark}
\newtheorem*{example}{Example}
\newcommand{\R}{\mathbb R}
\newcommand{\C}{\mathbb C}
\newcommand{\Z}{\mathbb Z}
\newcommand{\Q}{\mathbb Q}
\newcommand{\N}{\mathbb N}
\newcommand{\F}{\mathbb F}
\newcommand{\D}{\mathbb D}
\newcommand{\bB}{\mathbb B}
\newcommand{\bC}{\mathbb C}
\newcommand{\bR}{\mathbb R}
\newcommand{\bZ}{\mathbb Z}
\newcommand{\bN}{\mathbb N}
\newcommand{\bQ}{\mathbb Q}
\newcommand{\bD}{\mathbb D}
\newcommand{\bF}{\mathbb F}
\newcommand{\CP}{\mathbb{CP}}
\newcommand{\RP}{\mathbb{RP}}
\newcommand{\sO}{\mathscr O}
\newcommand{\sI}{\mathscr I}
\newcommand{\sF}{\mathscr F}
\newcommand{\sG}{\mathscr G}
\newcommand{\sH}{\mathscr H}
\newcommand{\sR}{\mathscr R}
\newcommand{\sM}{\mathscr M}
\newcommand{\sV}{\mathscr V}
\newcommand{\sU}{\mathscr U}
\DeclareMathOperator{\Hom}{Hom}
\DeclareMathOperator{\Ext}{Ext}
\DeclareMathOperator{\Tor}{Tor}
\DeclareMathOperator{\Spec}{Spec}
\DeclareMathOperator{\Proj}{Proj}
\DeclareMathOperator{\supp}{supp}
\DeclareMathOperator{\im}{im}
\DeclareMathOperator{\id}{id}
\DeclareMathOperator{\Tr}{Tr}
\DeclareMathOperator{\tr}{tr}
\DeclareMathOperator{\rank}{rank}
\DeclareMathOperator{\ord}{ord}
\DeclareMathOperator{\dist}{dist}
\DeclareMathOperator{\End}{End}
\DeclareMathOperator{\Aut}{Aut}
\DeclareMathOperator{\Gal}{Gal}
\DeclareMathOperator{\vol}{vol}
\DeclareMathOperator{\Cl}{Cl}
\DeclareMathOperator{\coker}{coker}
\DeclareMathOperator{\codim}{codim}
\DeclareMathOperator{\divergence}{div}
\DeclareMathOperator{\Ric}{Ric}
\DeclareMathOperator{\grad}{grad}
\DeclareMathOperator{\Hess}{Hess}
\newcommand{\abs}[1]{\left\lvert#1\right\rvert}
\newcommand{\sabs}[1]{\lvert#1\rvert}
\newcommand{\babs}[1]{\bigl\lvert#1\bigr\rvert}
\newcommand{\Babs}[1]{\Bigl\lvert#1\Bigr\rvert}
\newcommand{\bbabs}[1]{\biggl\lvert#1\biggr\rvert}
\newcommand{\BBabs}[1]{\Biggl\lvert#1\Biggr\rvert}
\newcommand{\norm}[1]{\left\lVert#1\right\rVert}
\newcommand{\snorm}[1]{\lVert#1\rVert}
\newcommand{\bnorm}[1]{\bigl\lVert#1\bigr\rVert}
\newcommand{\Bnorm}[1]{\Bigl\lVert#1\Bigr\rVert}
\newcommand{\bbnorm}[1]{\biggl\lVert#1\biggr\rVert}
\newcommand{\BBnorm}[1]{\Biggl\lVert#1\Biggr\rVert}
\newcommand{\widebar}[1]{\overline{#1}}
\newcommand{\nicefrac}[2]{\frac{#1}{#2}}
\newcommand{\linnprod}[2]{\langle#1,#2\rangle}
\newcommand{\myindex}[1]{#1}
\newcommand{\glsadd}[1]{}
\newcommand{\avoidbreak}{}
\newcommand{\sourceref}[1]{\textnormal{[原文：\nolinkurl{#1}]}}
\newcommand{\thmref}[1]{\sourceref{#1}}
\newcommand{\lemmaref}[1]{\sourceref{#1}}
\newcommand{\propref}[1]{\sourceref{#1}}
\newcommand{\corref}[1]{\sourceref{#1}}
\newcommand{\defnref}[1]{\sourceref{#1}}
\newcommand{\exampleref}[1]{\sourceref{#1}}
\newcommand{\exerciseref}[1]{\sourceref{#1}}
\newcommand{\figureref}[1]{\sourceref{#1}}
\newcommand{\sectionref}[1]{\sourceref{#1}}
\newcommand{\appendixref}[1]{\sourceref{#1}}
\newcommand{\stacksref}[2]{\href{https://stacks.math.columbia.edu/tag/#1}{#2 [#1]}}


\newcommand{\E}{\mathbb E}
\newcommand{\Prob}{\mathbb P}
\newcommand{\Reff}{\mathcal R_{\mathrm{eff}}}
\newcommand{\eps}{\varepsilon}
\DeclareMathOperator{\diam}{diam}
\DeclareMathOperator{\Var}{Var}
\DeclareMathOperator{\Crit}{Crit}
\newcommand{\dd}{\,\mathrm d}


\begin{document}
\section*{076\quad 次Gaussian独立坐标二次型的Hanson--Wright集中界}
\subsection*{来源与整理范围}
Mark Rudelson and Roman Vershynin, \emph{Hanson--Wright inequality and sub-gaussian concentration}, Electron. Commun. Probab. 18 (2013), no.82, DOI 10.1214/ECP.v18-2865。Theorem 1.1 与完整五步证明，印刷页2--5；记号取第1页。
\par\url{https://www.maths.tcd.ie/EMIS/journals/EJP-ECP/article/download/2865/2865-15071-1-PB.pdf}
\par\textbf{恢复方式（整理者说明）：}依据人类原文重新整理以补齐缺失题号；并非旧题证文件的逐字节恢复；2026-09-28。
\subsection*{作者命题与人类证明}

\paragraph{记号（据作者第1节整理）。}
$\|\xi\|_{\psi_2}=\sup_{p\geq1}p^{-1/2}(\mathbb E|\xi|^p)^{1/p}$，
$\|\xi\|_{\psi_1}=\sup_{p\geq1}p^{-1}(\mathbb E|\xi|^p)^{1/p}$。
$\|A\|$ 为 Euclidean 算子范数，$\|A\|_{\mathrm{HS}}=(\sum_{i,j}|a_{ij}|^2)^{1/2}$。
以下 $C,C_1,c,\ldots$ 为可变的正绝对常数。
\begin{theorem}[Hanson--Wright inequality; Theorem 1.1]
Let $X=(X_1,\ldots,X_n)\in\mathbb R^n$ be a random vector with independent
components $X_i$ which satisfy $\mathbb E X_i=0$ and $\|X_i\|_{\psi_2}\leq K$.
Let $A$ be an $n\times n$ matrix. Then, for every $t\geq0$,
\[
\mathbb P\{|X^{\mathsf T}AX-\mathbb E X^{\mathsf T}AX|>t\}
\leq2\exp\left[-c\min\left(\frac{t^2}{K^4\|A\|_{\mathrm{HS}}^2},
\frac{t}{K^2\|A\|}\right)\right].
\]
\end{theorem}
\begin{proof}
By replacing $X$ with $X/K$ we can assume without loss of generality that
$K=1$. Let us first estimate
\[
p:=\mathbb P\{X^{\mathsf T}AX-\mathbb E X^{\mathsf T}AX>t\}.
\]
Let $A=(a_{ij})_{i,j=1}^n$. By independence and zero mean of $X_i$, we can represent
\[
X^{\mathsf T}AX-\mathbb E X^{\mathsf T}AX
=\sum_i a_{ii}(X_i^2-\mathbb E X_i^2)+\sum_{i\ne j}a_{ij}X_iX_j.
\]
The problem reduces to estimating the diagonal and off-diagonal sums:
\[
p\leq\mathbb P\!\left\{\sum_i a_{ii}(X_i^2-\mathbb E X_i^2)>t/2\right\}
+\mathbb P\!\left\{\sum_{i\ne j}a_{ij}X_iX_j>t/2\right\}=:p_1+p_2.
\]
\textbf{Step 1: diagonal sum.}
Note that $X_i^2-\mathbb E X_i^2$ are independent mean-zero sub-exponential
random variables, and
\[
\|X_i^2-\mathbb E X_i^2\|_{\psi_1}
\leq2\|X_i^2\|_{\psi_1}\leq4\|X_i\|_{\psi_2}^2\leq4.
\]
These standard bounds can be found in the author's reference [18],
Remark 5.18 and Lemma 5.14. Then we can use a Bernstein-type inequality
([18], Proposition 5.16) and obtain
\begin{equation}\label{hw-diag}
p_1\leq\exp\!\left[-c\min\left(\frac{t^2}{\sum_i a_{ii}^2},
\frac{t}{\max_i|a_{ii}|}\right)\right]
\leq\exp\!\left[-c\min\left(\frac{t^2}{\|A\|_{\mathrm{HS}}^2},
\frac{t}{\|A\|}\right)\right].
\end{equation}
\textbf{Step 2: decoupling.}
It remains to bound the off-diagonal sum $S:=\sum_{i\ne j}a_{ij}X_iX_j$.
The argument will be based on estimating the moment generating function of $S$
by decoupling and reduction to normal random variables.
Let $\lambda>0$ be a parameter whose value we will determine later. By
Chebyshev's inequality,
\begin{equation}\label{hw-markov}
p_2=\mathbb P\{S>t/2\}\leq\exp(-\lambda t/2)\mathbb E\exp(\lambda S).
\end{equation}
Consider independent Bernoulli random variables $\delta_i\in\{0,1\}$ with
$\mathbb E\delta_i=1/2$. Since $\mathbb E\delta_i(1-\delta_j)$ equals $1/4$
for $i\ne j$ and $0$ for $i=j$, we have
\[
S=4\mathbb E_\delta S_\delta,\qquad
S_\delta=\sum_{i,j}\delta_i(1-\delta_j)a_{ij}X_iX_j.
\]
Here $\mathbb E_\delta$ denotes expectation with respect to
$\delta=(\delta_1,\ldots,\delta_n)$. Jensen's inequality yields
\begin{equation}\label{hw-jensen}
\mathbb E\exp(\lambda S)\leq\mathbb E_{X,\delta}\exp(4\lambda S_\delta).
\end{equation}
Consider $\Lambda_\delta=\{i\in[n]:\delta_i=1\}$ and express
\[
S_\delta=\sum_{i\in\Lambda_\delta,\ j\in\Lambda_\delta^c}a_{ij}X_iX_j
=\sum_{j\in\Lambda_\delta^c}X_j\left(\sum_{i\in\Lambda_\delta}a_{ij}X_i\right).
\]
Now condition on $\delta$ and $(X_i)_{i\in\Lambda_\delta}$. Then $S_\delta$
is a linear combination of mean-zero sub-gaussian random variables
$X_j$, $j\in\Lambda_\delta^c$, with fixed coefficients. Its conditional
sub-gaussian norm is bounded by the $\ell^2$ norm of the coefficient vector:
\[
\|S_\delta\|_{\psi_2}\leq C\sigma_\delta,\qquad
\sigma_\delta^2:=\sum_{j\in\Lambda_\delta^c}
\left(\sum_{i\in\Lambda_\delta}a_{ij}X_i\right)^2.
\]
The standard moment generating function estimate for centered sub-gaussians
([18], Lemma 5.5) yields
\[
\mathbb E_{(X_j)_{j\in\Lambda_\delta^c}}\exp(4\lambda S_\delta)
\leq\exp(C\lambda^2\|S_\delta\|_{\psi_2}^2)
\leq\exp(C'\lambda^2\sigma_\delta^2).
\]
Taking expectations with respect to $(X_i)_{i\in\Lambda_\delta}$, we obtain
\begin{equation}\label{hw-Ed}
\mathbb E_X\exp(4\lambda S_\delta)
\leq\mathbb E_X\exp(C'\lambda^2\sigma_\delta^2)=:E_\delta.
\end{equation}
This holds for every fixed $\delta$. It remains to estimate $E_\delta$.

\textbf{Step 3: reduction to normal random variables.}
Consider $g=(g_1,\ldots,g_n)$ where $g_i$ are independent $N(0,1)$ variables.
For each fixed $\delta$ and $X$, rotation invariance gives
\[
Z:=\sum_{j\in\Lambda_\delta^c}g_j
\left(\sum_{i\in\Lambda_\delta}a_{ij}X_i\right)
\sim N(0,\sigma_\delta^2).
\]
Thus $\mathbb E_g\exp(sZ)=\exp(s^2\sigma_\delta^2/2)$.
Choosing $s^2=2C'\lambda^2$ and comparing with \eqref{hw-Ed} gives
\[
E_\delta=\mathbb E_{X,g}\exp(C_1\lambda Z),\qquad C_1=\sqrt{2C'}.
\]
Rearrange $Z=\sum_{i\in\Lambda_\delta}X_i
(\sum_{j\in\Lambda_\delta^c}a_{ij}g_j)$.
Bound its moment generating function as in Step 2, now using the
sub-gaussian properties of the $X_i$ for $i\in\Lambda_\delta$. We obtain
\[
E_\delta\leq\mathbb E_g\exp\left[C_2\lambda^2
\sum_{i\in\Lambda_\delta}\left(\sum_{j\in\Lambda_\delta^c}a_{ij}g_j\right)^2\right].
\]
Let $P_\delta$ be the coordinate projection onto $\mathbb R^{\Lambda_\delta}$
and $A_\delta=P_\delta A(I-P_\delta)$. Then
\[
E_\delta\leq\mathbb E_g\exp(C_2\lambda^2\|A_\delta g\|_2^2).
\]
The original variables $X_i$ have now been removed from the problem.

\textbf{Step 4: calculation for normal random variables.}
By rotation invariance, $\|A_\delta g\|_2^2$ is distributed identically
with $\sum_i s_i^2g_i^2$, where $s_i$ are the singular values of $A_\delta$.
By independence,
\[
E_\delta\leq\mathbb E_g\exp\left(C_2\lambda^2\sum_i s_i^2g_i^2\right)
=\prod_i\mathbb E\exp(C_2\lambda^2s_i^2g_i^2).
\]
Each $g_i^2$ is chi-squared with one degree of freedom, so
$\mathbb E\exp(tg_i^2)=(1-2t)^{-1/2}$ for $t<1/2$. Therefore
\[
E_\delta\leq\prod_i(1-2C_2\lambda^2s_i^2)^{-1/2}
\quad\hbox{provided }\max_i C_2\lambda^2s_i^2<1/2.
\]
Using $(1-z)^{-1/2}\leq e^z$ for $0\leq z\leq1/2$, this becomes
\[
E_\delta\leq\prod_i\exp(C_3\lambda^2s_i^2)
=\exp\left(C_3\lambda^2\sum_i s_i^2\right)
\quad\hbox{provided }\max_i C_3\lambda^2s_i^2<1/2.
\]
Since $\max_i s_i=\|A_\delta\|\leq\|A\|$ and
$\sum_i s_i^2=\|A_\delta\|_{\mathrm{HS}}^2\leq\|A\|_{\mathrm{HS}}^2$,
\[
E_\delta\leq\exp(C_3\lambda^2\|A\|_{\mathrm{HS}}^2)
\quad\hbox{for }\lambda\leq c_0/\|A\|.
\]
This is uniform in $\delta$. Taking expectations and recalling
\eqref{hw-jensen} and \eqref{hw-Ed}, we obtain
\[
\mathbb E\exp(\lambda S)\leq\mathbb E_\delta E_\delta
\leq\exp(C_3\lambda^2\|A\|_{\mathrm{HS}}^2),\qquad \lambda\leq c_0/\|A\|.
\]
\textbf{Step 5: conclusion.}
Putting this into \eqref{hw-markov} gives
\[
p_2\leq\exp(-\lambda t/2+C_3\lambda^2\|A\|_{\mathrm{HS}}^2),
\qquad\lambda\leq c_0/\|A\|.
\]
Optimizing over $\lambda$,
\[
p_2\leq\exp\left[-c\min\left(\frac{t^2}{\|A\|_{\mathrm{HS}}^2},
\frac{t}{\|A\|}\right)\right]=:p(A,t).
\]
Together with \eqref{hw-diag}, $p\leq p_1+p_2\leq2p(A,t)$.
Repeating for $-A$ gives the same bound for the lower tail. Combining them,
\[
\mathbb P\{|X^{\mathsf T}AX-\mathbb E X^{\mathsf T}AX|>t\}\leq4p(A,t).
\]
Finally one can reduce the factor $4$ to $2$ by adjusting $c$ in $p(A,t)$.
The proof is complete.
\end{proof}

\subsection*{整理说明与先修范围（非作者正文）}
允许引用独立次Gaussian和的范数界、中心化平方的次指数范数界、Bernstein不等式与Gaussian矩母函数；不把解耦、Gaussian替换或奇异值乘积估计藏进先修。本题需独立完成这些构造及参数优化，定位为高级概率的综合证明。保留作者五步次序；省略与主证明无关的历史综述和第二节应用。原PDF式(1.1)首项漏印$\exp$、Step 4第一式将已得上界印为等号、第5页首行末项漏平方、末段$p=p_1+p_2$应为$\leq$，此转录按同段已定义量修正并在此明列。参考文献[18]为Roman Vershynin, \emph{Introduction to the non-asymptotic analysis of random matrices}, 2012。
\subsection*{可修改的简短标注（非作者正文）}
$c$：随机二次型的尾概率控制

$a$：Bernoulli随机分割；条件矩母函数；Gaussian替换；奇异值分解；Chernoff优化

\texttt{cross\_theory\_status = yes}。作者Steps 2--4把相关的二阶乘积和转为独立两组变量的双线性型，再转为Gaussian奇异值乘积，使谱范数与Hilbert--Schmidt范数共同控制原尾概率。

全部标注为非穷尽、可修改初稿。
\subsection*{来源许可与编辑归属}
原期刊文章页明确标示CC BY 3.0：\url{https://emis.muni.cz/journals/EJP-ECP/article/view/2865.html}；许可\url{https://creativecommons.org/licenses/by/3.0/}。
ChatGPT 加入题号、中文来源说明、整理范围和初稿标注；编辑说明与人类证明分列。
\end{document}
